\section{Motivation and the gap}

We want to predict which residues of an antigen are B-cell epitope residues, from the antigen
structure alone, in cases where the antibody is unknown. Doing this rigorously requires
residue-level labels for ``binding importance.'' We take these from \textbf{alanine scanning}:
for residue $i$, $\Delta\Delta G_\mathrm{bind}^{(i)} = \Delta G_\mathrm{bind}^{\mathrm{mut},i} -
\Delta G_\mathrm{bind}^{\mathrm{WT}}$, with large positive values flagging residues whose
wild-type identity stabilizes the interface. We treat this as a residue-level epitope score,
$S_i^\mathrm{epitope} \approx \Delta\Delta G_\mathrm{bind}^{(i)}$. (Full derivation and the
alchemical / StaB-ddG label tiers: \texttt{../../tcr/docs/b\_cell\_eptiope\_prediction\_may23\_update\_final.pdf}.)

\paragraph{The gap.} Most B-cell epitope predictors label residues by \emph{contact distance} to
a bound antibody. That conflates ``close'' with ``important'': a residue can sit at the interface
yet contribute little energetically, or be energetically critical without being the nearest
neighbor. Alanine-scan $\Delta\Delta G$ is a cleaner, energetic definition of importance. Separately,
since Westhof et al.~\cite{westhof1984} it has been known that antigenicity correlates with
\emph{segmental mobility} (flexibility / B-factor) --- an antigen-intrinsic, antibody-agnostic
property. The opportunity is to test, with modern MD, whether antigen-only \emph{dynamical} and
\emph{solvation} features (RMSF, collective-mode participation, hydration thermodynamics) predict
energetic binding importance --- and whether they beat static exposure. This is where the
author's methods background (FRESEAN collective-mode analysis; 3D-2PT-style hydration
thermodynamics; non-equilibrium energy flow) is the differentiator
(\texttt{../../\_literature/my\_manuscripts/}).
